\documentclass[instructions.tex]{subfiles}
\begin{document}

\chapter{Motifs de conversion, tirés de la passion du Sauveur.}

\section{Ce qu’il a souffert en son âme.}

Il a fallu, dit saint Paul, que Jésus-Christ souffrît ; il l’a fallu pour nous apprendre ce que c’est que Dieu, ce que c’est que l’homme, ce que c’est que le péché. Non, sans les souffrances de Jésus-Christ, nous n’aurions jamais pu connoître quelle est la grandeur et la sainteté de Dieu, quelle est la dignité et en même temps la misère de l’homme, quelle est la noirceur et la malice du péché. Il a donc fallu qu’il souffrît\footnote{Christum oportuit pati. Act. 17. 3.}.

Mais qu’a-t-il souffert ? Il a souffert tout ce qu’il a pu souffrir en son âme et en son corps. Dans son âme, il a ressenti l’horreur la plus insupportable, la douleur la plus profonde, la frayeur la plus excessive, la confusion la plus humiliante, et tout cela au-dessus de ce que l’esprit humain peut comprendre.

\primo L’horreur. S’étant offert à Dieu son Père pour être notre caution, il se voyoit chargé de nos dettes, comme investi et revêtu de nos péchés. Quels seroient les sentiments d’un jeune prince qu’on revêtirait d’habits sales, remplis d’ordure et d’infection, si on l’obligeoit à les porter toute sa vie, et à être enfermé dans un lieu où il ne verroit que des monstres et des objets hideux ? Quelle horreur n’auroit-il pas de se voir en cet état ! Le Fils de Dieu a eu infiniment plus d’horreur de se voir chargé des péchés de l’univers ; il étoit comme plongé et noyé dans les ordures et dans les iniquités de tous les hommes, qui sans cesse se présentoient à son esprit, sans que la vue de ces affreux objets l’ait quitté un seul moment\footnote{Infixus sum in limo profundi ; tempestas demersit me. Ps. 68.}.

\secundo La douleur. Tous nos péchés réunis, dans sa personne, selon l’expression du Prophète, le pénétrèrent d’une douleur si vive, que la vue d’un seul péché et d’une âme perdue, étoit capable de lui causer la mort. Oh ! quelle fut donc son affliction, quelles furent ses larmes, à la vue de la multitude immense des péchés de tous les hommes et de tous les siècles, à la vue de la perte de tant d’âmes qui, par leur malice, devoient abuser de son sang\footnote{Quae utilitas in sanguine meo ? Ps. 29.} ! Il en fut accablé d’une tristesse si grande, qu’il tomba dans l’agonie, et en sua le sang. Sa douleur fut si profonde, que, si elle avoit été partagée entre tous les hommes, aucun n’auroit pu la soutenir ; tous en seroient morts.

\tertio La frayeur. Jésus-Christ avoit continuellement présentes à son esprit la justice et la sainteté de son Père offensé, les châtimens que méritoient nos offenses, et les tourmens horribles qu’il devoit souffrir pour les expier. Ces pensées affligeantes accabloient intérieurement son âme pendant tout le cours de sa vie ; mais, au jardin des Oliviers, il permit que tous ces objets fissent une impression si vive dans son imagination, que son humanité sainte en succomba, qu’il en trembla lui-même\footnote{Coepit pavere et taedere. Marc. 14. 33.} ; sans miracle, l’accablement et la frayeur l’auroient fait expirer dans le moment.

\quarto Mais ce qui est le plus incompréhensible dans cet Homme-Dieu, c’est la confusion qu’il a ressentie\footnote{Operuit confusio faciem meam. Ps. 2.}. \primo du côté de Dieu son Père, dont la sainteté l’accabloit tellement de confusion, en se voyant chargé des iniquités du monde, qu’il se considéroit, devant son Père, aussi criminel que s’il eût commis tous les péchés des hommes, se regardant comme un objet d’anathème et de malédiction, dit saint Paul\footnote{Factus pro nobis maledictum. Gal. 3.}. \secundo Confusion du côté de lui-même. Un homme qui seroit couvert d’une lèpre hideuse, seroit saisi d’une telle honte, qu’il n’oseroit même s’exposer aux yeux du public. Oh ! combien grande fut la confusion intérieure du Fils de Dieu, se voyant couvert de nos souillures, et absorbé dans la masse de nos crimes ! objets mille fois plus humilians que la lèpre la plus honteuse. Saint Paul a donc bien raison de dire que Jésus-Christ n’avoit aucune complaisance en lui-même, parce que les crimes et les outrages des hommes étoient tombés sur lui\footnote{Christus non sibi placuit, sed, sicut scriptum est, improperia improperantium tibi ceciderunt super me. Rom. 15.}.

\tertio Enfin, du côté des hommes, quelle fut sa confusion de voir sa sainteté et sa divinité méprisées, et sa personne innocente et adorable accusée de crimes, devenue le jouet d’une populace effrénée ; de se voir bafoué, insulté, moqué, traîné par les rues, lié comme un malfaiteur et un scélérat ; foulé aux pieds comme un ver de terre ; crucifié comme un infâme, comme l’opprobre des humains !

O hommes, que pensez-vous à la vue d’un Dieu affligé et traité avec tant d’ignominie ? Voilà l’état où il est réduit pour expier vos péchés, et pour vous en faire comprendre l’énormité ; pouvez-vous, après cela, compter le péché pour rien ?

\section{Résolutions}

\primo Je méditerai de temps en temps sur les sentimens d’horreur, de douleur, de frayeur et de confusion, dont l’âme du Sauveur fut si tourmentée dans sa passion, à cause de mes péchés.  

\secundo Et je relirai ce chapitre chaque fois que je voudrai aller à confesse, pour m’approprier ces mêmes sentimens. 

\prayer{O Jésus ! faites-moi partager d’une manière salutaire les douleurs intérieures que vous ressentîtes dans la carrière de vos pénibles souffrances, lorsque vous vous considériez comme chargé du fardeau hideux de toutes mes iniquités.}
\end{document}
